\begin{frame}{이 숫자의 의미는 비교해야 보인다}
  \begin{itemize}
    \item 같은 8$\times$8을 \textbf{flatten(64-dim)}해서
      64$\to$64$\to$10 단순 MLP면 파라미터 \textbf{4,810}개뿐 —
      conv 층(320)보다 \textbf{작다}
    \item 즉, 작은 데이터에서는 ``깊게/넓게''가 아니라
      \textbf{입력 차원(flatten)이 작으면 GBDT}(특수한 파라미터가
      아예 없음)가 효율적이다
    \item ``왜 CNN''의 진짜 이유 = \textbf{해상도 $S$가 커질 때}
  \end{itemize}
  \vskip0.3em
  \small
  \begin{tabular}{@{}cccc@{}}
    \toprule
    해상도 $S$ & ``flatten $\to$ 256'' FC ($S^2\times256$) &
      3$\times$3 conv(1$\to$32) & 비율 \\
    \midrule
    8 & 16,384 & 320 & 51$\times$ \\
    32 & 262,144 & 320 & 819$\times$ \\
    64 & 1,048,576 & 320 & 3,277$\times$ \\
    224 & 12,845,056 & 320 & 40,141$\times$ \\
    \bottomrule
  \end{tabular}
  \vskip0.5em
  \begin{block}{``상수 vs 제곱''}
    flatten은 해상도 \textbf{제곱}($S^2$)으로 폭증하지만, conv는
    해상도와 무관하게 \textbf{320으로 일정} — 이 차이가
    ``이미지가 커질수록 CNN(파라미터 공유, 평행불변)이 필수''를
    \textbf{수학적으로} 말한다. 이 표 한 장이 발표용
    \textbf{수식적 정당화}로 그대로 쓰인다.
  \end{block}
\end{frame}
